\section{材料与方法}
\label{sec:methods}

\subsection{数据与分析窗口}
\label{sec:meth-data}

每周实验室确诊流感住院人数及报告医院百分比取自美国国家医疗安全网络（NHSN）的 Weekly Hospital Respiratory Data 数据表，覆盖美国 50 个州和哥伦比亚特区\citep{cdc2026nhsn}，所用版本于 2026-10-20 获取。数据表共含 2020-08-08 至 2026-10-10 的 21641 个辖区周。2024 年 4 月 30 日前为强制报告期，2024 年 5 月 1 日至 10 月 31 日为自愿报告期，2024 年 11 月 1 日起恢复强制报告；在三个报告期内，报告规模不低于五的辖区周分别为 6009、420 和 4077 个。

主分析窗口覆盖 2022--23 至 2025--26 四个流行季的流行早期增长阶段。每个窗口自 10 月 1 日起（2024 年则自强制报告恢复的 11 月 9 日起），包含全国周增长比不低于 1.05 的连续周，并于全国首个转折点的前一周结束。各窗口依次为 2022-10-01 至 2022-11-26、2023-10-07 至 2023-12-23、2024-11-09 至 2024-12-28 以及 2025-10-04 至 2025-12-27，共 42 个日历周。上述规则取代了最初版本；最初版本以流行季最大值作为各窗口终点，在双峰的 2024--25 季仅得到四周。规则于 2026-10-20 修订，修订时尚未在新窗口上拟合任何截距。为保持连续性，本分析早期版本的三个日历窗口（2022-10-08 至 2022-11-26、2023-10-14 至 2023-12-16 以及 2024-11-09 至 2025-01-18）作为次要设定进行分析。峰值窗口包含各流行季的全国峰值周及其后七周。所有窗口均未纳入以零填补的病例数。

\subsection{周增长统计量}
\label{sec:meth-stats}

对每个州和每一周计算两个周增长统计量。加权增长比定义为：
\begin{equation}
  G_t=\frac{C_{t+1}}{D_t},\qquad D_t=0.65\,C_t+0.25\,C_{t-1}+0.10\,C_{t-2}, % qa-allow
  \label{eq:dt}
\end{equation}
其中 $C_t$ 为第 $t$ 周的住院人数，$D_t$ 为报告规模，且 $\sum_sw_s{}^2=0.495$。$G$ 仅使用 $D_t\ge5$ 的周。更新型比值借助代际间隔将单周增长换算为再生数\citep{wallinga2007how}，其定义为：
\begin{equation}
  R_{A,t}=\left(1+r_t\theta\right)^{\alpha},\qquad r_t=\frac{1}{7}\log\frac{C_{t+1}+0.5}{C_t+0.5}, % qa-allow
  \label{eq:ra}
\end{equation}
其中 $r_t$ 为日增长率；代际间隔服从均值 3.6 天、标准差 1.6 天的伽马分布（取 \citet{cowling2009estimation} 估计的家庭序列间隔），$\alpha$ 与 $\theta$ 分别为其形状参数与尺度参数。敏感性分析采用均值为 2.2 天和 4.3 天的代际间隔。$R_A$ 以单周计数为规模，因此只在 $D_t\ge5$ 且 $C_t\ge5$ 的周计算。$R_A$ 由观测增长变换得到，并非 $R_t$ 的后验估计\citep{cori2013new}；两种统计量因此都属于表观量，它们与 $v_R$ 的关系由模拟确定（第~\ref{sec:meth-inversion}~节）。

\subsection{期望损失与平衡规模}
\label{sec:meth-loss}

常规监测中原发病例数无法观测，继发感染来源于感染压力 $\Lambda_{t+1}=\sum_{s\ge1}w_sI_{t+1-s}$\citep{fraser2007estimating, cori2013new}。给定历史 $\mathcal{F}_t$，对 $C_{t+1}$ 应用全方差公式，并令 $m=\rho\E(\Lambda_{t+1}\mid\mathcal{F}_t)$，相对预测方差可写为：
\begin{equation}
  \frac{\Var(C_{t+1}\mid\mathcal{F}_t)}{m^2}=v_R+\frac{A}{m}+\mathrm{CV}^2(\Lambda_{t+1}\mid\mathcal{F}_t)\,(\Rbar^2+v_R),
  \label{eq:info-decomp}
\end{equation}
其中 $\mathrm{CV}$ 为变异系数。最后一项是另一个与尺度无关的分量，仅在感染压力已知时消失。对于均值为 $\bar{X}$ 的统计量 $X_t$，相对损失 $L_t=[(X_t-\bar{X})/\bar{X}]^2$ 的模型设定为：
\begin{equation}
  \E(L_t\mid D_t)=a+\frac{b}{D_t},
  \label{eq:affine}
\end{equation}
其中，若一周恰为一代且 $D_t$ 不含噪声，则 $a=v_R/\Rbar^2$、$b=A/\Rbar^2$，且 $b/a=\mx$。分母中二项噪声的相对方差为 $(1-\rho)\sum_sw_s{}^2/m$，这一噪声使斜率近似增大为：
\begin{equation}
  b\approx\frac{A+(1-\rho)(\Rbar^2+v_R)\sum_sw_s{}^2}{\Rbar^2}.
  \label{eq:b-delta}
\end{equation}
计算 $G$ 的 $b/a$ 理论范围时，$\Rbar$ 取各设定下 $G$ 的均值，$v_R\in\{0.04,0.08\}$，$k\in\{0.2,0.5\}$\citep{lloydsmith2005superspreading}，每例感染的住院概率 $\rho\in\{0.0056,0.021\}$\citep{tokars2018seasonal}，并分别计算使用与不使用式~(\ref{eq:b-delta}) 的结果。在这些情景下，未经校正的平衡规模为 21.3--49.2 例。

\subsection{估计与推断}
\label{sec:meth-inference}

式~(\ref{eq:affine}) 采用可行广义最小二乘拟合，权重为 $1/\hat{\mu}^2$，其中 $\hat{\mu}$ 为上一轮迭代的拟合均值，迭代直至收敛\citep{wooldridge2010econometric}。同一州的观测在时间上相关，同一周的观测共享全国性影响。因此，不确定性采用州 $\times$ 周双向自助法估计，即对州和周分别独立有放回抽样，并以两次抽取次数之积作为每个观测的权重\citep{cameron2008bootstrap, cameron2011robust}。每个设定使用 2000 次自助重复，并各用独立的随机种子（规则窗口中 $R_A$ 为 20261024、$G$ 为 20261124，日历窗口中 $G$ 为 20261224、$R_A$ 为 20261324）；不确定性报告为点估计及其 95\% 百分位区间，例如 0.0150（95\% CI: 0.0100, 0.0208）。

截距分解通过两次重新拟合式~(\ref{eq:affine}) 完成。第一次重新拟合使用各周内相对跨州均值的偏差，以去除全国周效应，并乘以小样本因子 $n_w/(n_w-1)$。第二次重新拟合去除以交替投影估计的可加周效应与州-季效应。全国周分量为未调整截距与周调整截距之差，州-季分量为加入州-季效应后截距的进一步变化。仿射模型与幂律 $\E(L_t\mid D_t)=cD_t{}^{-\gamma}$ 的比较采用两种方法：一是残差平方和的加权对比，并以幂律下的残差自助法校准；二是逐州留一与逐季留一交叉验证。检测正截距的统计功效由拟合的仿射模型模拟得到，模拟所用截距取不超过估计值的各个取值。$b/a$ 的敏感性分析包括剔除影响最大的 1\% 损失值、每次剔除一个流行季、使用 Huber 权重以及去除全国周效应。补充分析使用峰值窗口，改变代际间隔，并将分母权重改为单周或等权三周。

\subsection{报告完整度}
\label{sec:meth-reporting}

将各州各周的病例数除以对应的报告医院比例后，计算覆盖校正统计量。截距变化 $\Delta a=a-a_{\mathrm{adj}}$ 与零的等效性采用 5\% 水平的双单侧检验（TOST）评估，所用 90\% 区间来自配对双向自助法\citep{schuirmann1987comparison}。等效界取为按命题~\ref{prop:report-noise} 相对报告波动 $\tau_0$ 所增加的截距，具体为：
\begin{equation}
  \delta_G=\tau_0{}^2\left(1+\frac{v_R}{\Rbar^2}\right),\qquad % qa-allow
  \delta_{R_A}=\kappa_A{}^2\,\tau_0{}^2\left(1+\frac{v_R}{\Rbar^2}\right),
  \label{eq:tost}
\end{equation}
其中 $\kappa_A=\mathrm{d}\log R_A/\mathrm{d}\log g$ 将周增长 $g$ 的扰动传递至 $R_A$，$\Rbar$ 为各统计量的样本均值，$v_R=0.04$，$\tau_0=0.10$（主分析）或 0.05。由于 $\Rbar$ 为估计值，以截距份额表示的等效界也带有自助法区间。功效采用正态近似计算。比较 2023 年、2024 年与 2025 年 5 月至 10 月间 $\log G$ 的四分位距与 $G$ 的标准差。作为第二个代理指标，利用 44 个有数据的州的临床实验室每周流感检测阳性数计算 $R_A$；以医院 $R_A$ 对实验室 $R_A$ 回归所去除的医院截距比例，采用 500 次双向自助重复估计。

\subsection{流行季与病原体}
\label{sec:meth-seasons}

各流行季截距采用 DerSimonian--Laird 随机效应模型合并\citep{dersimonian1986meta}；合并截距采用自由度为 $S-1$ 的 $t$ 区间，季间方差采用 Q-profile 区间。推论~\ref{cor:season} 给出任一季间方差估计相对标准误的下界\citep{searle1992variance}。截距比 $q_1=a_{(1,0,0)}/a_{w}$ 与 $q_2=a_{(1/3,1/3,1/3)}/a_{w}$ 分别将单周分母与等权三周分母同式~(\ref{eq:dt}) 的权重相比较。流感在四个窗口中计算这两个截距比，COVID-19 与 RSV 住院数据各在两个窗口中计算。零分布范围来自含报告噪声、无环境变异的流行增长模拟，模拟遍历三种代际间隔、三种住院延迟和四个噪声水平。八个组合各按五项单侧准则检验（$q_1<1$、$q_2>1$、两个比值各自落在零分布范围之外，以及相对 $v_R/\Rbar^2$ 的正超额），共 40 项检验，采用 Holm 与 Benjamini--Hochberg 程序校正\citep{holm1979simple, benjamini1995controlling}。

\subsection{表观统计量的模拟}
\label{sec:meth-sim}

第~\ref{sec:background}~节的闭式结果通过蒙特卡洛模拟检验，每个点重复 20000 次。在第一套通用设计中，表观周统计量由日尺度更新模型模拟生成，取 $\Rbar=1.2$，代际间隔同上。每周 $R_t$ 在各周间独立地从方差为 $v_R$ 的伽马分布中抽取；后代服从负二项分布；住院经均值 5.5 天、标准差 3.0 天的伽马延迟后以概率 $\rho_h$ 二项发生，延迟分布的参数为假定值。每个重复包含 500 个州，初始报告规模在 2 与 150 之间服从对数均匀分布，另含四个预热周与六个观测周。十七个情景由 $v_R\in\{0.01,0.02,0.04,0.08,0.16\}$ 与 $k\in\{0.1,0.35,1\}$ 在 $\rho_h=0.01$ 下交叉组合，并在基线设定下增加 $\rho_h\in\{0.005,0.02\}$，每个情景重复 20 次。每个截距的无噪声极限来自 20000 个州的平均场运行，其中全部计数噪声均以期望值代替。

\subsection{两步反演}
\label{sec:meth-inversion}

实证截距分两步映射为 $v_R$。第一步，以数据匹配版本的模拟器给出各统计量期望拟合截距随 $v_R$ 变化的函数。每个模拟面板的横截面与长度和分析面板相同，包含 51 个州和四个规则窗口，观测周数依次为 9、12、8 和 13。流行季均值 $R_s$ 取各窗口全国周增长几何均值的更新型变换，各州初始住院数等于该州在窗口首周的观测住院数。周再生数为 $R_{i,s,w}=R_s+\eta_{s,w}+\varepsilon_{i,s,w}$，其中全国周冲击 $\eta$ 的方差为 $0.45\,v_R$，为所有州共有；州冲击 $\varepsilon$ 的方差为 $0.55\,v_R$；取值在 0.05 处截断。0.45 的份额取在两种统计量观测到的全国周份额 0.41 与 0.51 之间（第~\ref{sec:res-multiscale}~节），因此模拟复现数据的组内相关系数与全国周份额只是一项一致性检查；把该份额设为零或 0.65 时，$R_A$ 的隐含方差由 0.040 变为 0.039 和 0.040（表~\ref{tab:s8}）。传播为后代服从负二项分布（$k=0.35$）的日尺度更新过程，住院经上述延迟后以概率 $\rho_h=0.01$ 二项发生；两种统计量的计算与筛选完全按数据处理，分析者采用均值 3.6 天的代际间隔。设计取 $v_R\in\{0.02,0.04,0.06,0.09\}$，每个取值 200 个面板重复，对每个重复拟合式~(\ref{eq:affine})。校准线 $\E(\hat{a})=\gamma_0+\gamma_1v_R$ 采用加权最小二乘拟合，权重为各均值蒙特卡罗方差的倒数；尺度因子 $\kappa=\E(\hat{a})/(v_R/\Rbar^2)$ 中的 $\Rbar$ 取按观测行加权的真实均值。

第二步，对各统计量计算隐含方差 $v_R^{\mathrm{impl}}=(a-\gamma_0)/\gamma_1$。其区间合并了两部分变异：两个实证截距的 2000 次联合双向自助抽样（两者共用州与周的抽取次数），以及 $(\gamma_0,\gamma_1)$ 按其蒙特卡罗抽样分布的独立抽取。两种统计量的一致定义为未拒绝，即两者之差的 95\% 区间包含零。对主设计另以双单侧检验检验等效：以 0.80 与 1.25 为隐含方差之比的界，使用比值抽样的 90\% 区间。敏感性设定每次只改变设计中的一个要素（表~\ref{tab:s8}），包括去除全国冲击或将其份额提高到 0.65，以及加入 0.5 的周间一阶自相关。分析者始终使用 3.6 天的代际间隔，真实代际间隔均值则设为 2.2 天或 4.3 天；延迟均值设为 3.5 天或 8 天。另有一个设定把 $k$ 降至 0.1，一个设定以作用于报告计数的同步乘性冲击代替传播冲击，上述通用设计则作为衔接设定。终修时新增四个设定：真实代际间隔取 \citet{cowling2009estimation} 区间下端的 2.9 天，一阶自相关取 0.25，以及每周住院数取自方差均值比为 1.5 或 2.0 的伽马--泊松混合（表~\ref{tab:s10}）。周冲击的独立性以各统计量 $x/\bar{x}-1$ 在州--季节块内的汇总滞后一阶自相关检验，区间来自 1000 次按州整群的自助抽样，并与 $v_R=0.04$ 时模拟面板中的分布比较。校准后的理论 $b/a$ 采用同一设计，取 $v_R\in\{0.04,0.08\}$、$k\in\{0.2,0.5\}$ 与 $\rho_h\in\{0.0056,0.021\}$；单个重复间 $b/a$ 的取值范围给出同等规模单个面板的预期分布。模拟种子由 20261020 经按设定、$v_R$ 取值与重复块索引的种子序列派生；联合自助法与校准抽取分别使用 20261921 和 20261922。
